\section{Additional Considerations, Abbreviations, and Amendment History}

\subsection{Additional Considerations}

Although this protocol is written to the United States combined IND/IDE pathway,
the LLM-directed robotic system is a Software as a Medical Device (SaMD) whose
deployment may be sought across jurisdictions, and the design anticipates the
parallel regulatory tracks summarized in Table~\ref{tab:jurisdictions}. The
European Union Medical Device Regulation (EU MDR) and its companion artificial-intelligence
requirements, China's National Medical Products Administration (NMPA), Japan's
Pharmaceuticals and Medical Devices Agency (PMDA), the United Kingdom's Medicines
and Healthcare products Regulatory Agency (MHRA), Health Canada, and Australia's
Therapeutic Goods Administration (TGA) each maintain a SaMD and robotically
assisted surgery framework with which the device evidence is designed to remain
compatible. The common technical spine, Phase 0 simulation validation, USL
readiness rating, hash-chained audit trail, and verification-before-generation
gate surface, is constructed to map onto each track without redesign, and the
international consensus standards cited in \S12 (IEC 80601-2-77, ASME V and V 40,
NASA-STD-7009A, IEEE 7009-2024) provide shared basis across jurisdictions.

The trial is additionally aligned with the legislative and financial framework of
the Verification Before Generation in Physical AI Oncology Trials Act of 2026 (H.
R. 9510) \cite{kawchak2026hr9510}, which ties the funding and oversight of Physical
AI oncology investigations to a verification-and-validation-and-uncertainty-quantification
(VVUQ) evidentiary standard. The ten-gate assurance suite operationalizes that
standard: every proposed robot-patient command is verified before it is generated,
against fourteen external safety and robotics standards plus two clinical
baselines, with epistemic and aleatory uncertainty bounds, and the gate surface
returns one of accept, block, or escalate. Figure~12 renders the ten-gate
pipeline.

\begin{table}[htbp]
\centering
\footnotesize
\caption{Cross-jurisdictional Software as a Medical Device (SaMD) and robotically
assisted surgery tracks. The common technical spine (Phase 0 validation, USL readiness, the
hash-chained audit trail, and the verification-before-generation gate surface)
maps onto each track without redesign.}
\label{tab:jurisdictions}
\begin{tabularx}{\textwidth}{L{2.5cm} L{3.2cm} Y}
\toprule
\textbf{Jurisdiction} & \textbf{Regulator / framework} & \textbf{Alignment basis} \\
\midrule
European Union & EU MDR + AI reqs & SaMD risk classification; common simulation-validation, audit evidence. \\
China & NMPA & SaMD, robotically assisted surgery review; standards-based VVUQ. \\
Japan & PMDA & SaMD and device review; consensus-standard alignment. \\
United Kingdom & MHRA & SaMD framework; post-market and audit-trail compatibility. \\
Canada & Health Canada & SaMD licensing; quality-system and validation evidence. \\
Australia & TGA & SaMD inclusion; standards-based safety and performance evidence. \\
\bottomrule
\end{tabularx}
\end{table}

\vspace{-0.2cm}

\subsection{Protocol Amendment History}

This is the initial version of the protocol. Table~\ref{tab:amend} records the
version history; subsequent amendments will be appended with the version, date,
sections affected, and a summary of the change.

\begin{table}[htbp]
\centering
\footnotesize
\caption{Protocol amendment history.}
\label{tab:amend}
\begin{tabularx}{\textwidth}{L{2.0cm} L{2.6cm} L{3.2cm} Y}
\toprule
\textbf{Version} & \textbf{Date} & \textbf{Sections affected} & \textbf{Summary of change} \\
\midrule
v1.0.0 & June 20, 2026 & All (initial release) & Initial version of the protocol. Establishes the combined IND/IDE design, the perioperative daraxonrasib 3+3 dose-finding arm, the on-premises LLM-directed eight-arm robotic Whipple device arm, the Physical AI Subpart J overlay, the statistical and oversight framework, and the back matter. \\
\bottomrule
\end{tabularx}
\end{table}

\begin{mermaidfig}
\node[mmin]  (gen)  at (-0.2,0)      {Proposed robot-patient\\interaction code /\\command};
\node[mmstep](gates)at (4.6,0)    {10-gate assurance suite\\14 external standards\\+ 2 clinical baselines};
\node[mmstep](uq)   at (9.6,0)    {Uncertainty\\quantification:\\epistemic + aleatory};
\node[mmdec] (v)    at (9.6,-2.6) {Gate surface\\verdict};
\node[mmgoal](acc)  at (4.6,-2.6) {ACCEPT: verified\\before generation};
\node[mmdark](blk)  at (-0.2,-2.6)   {BLOCK: command\\refused};
\node[mmdark](esc)  at (4.6,-4.5) {ESCALATE: hand\\back to human};
\draw[mmedge] (gen) -- (gates);
\draw[mmedge] (gates) -- (uq);
\draw[mmedge] (uq) -- (v);
\draw[mmedge] (v) -- node[above,font=\scriptsize, xshift=0.1cm]{accept} (acc);
\draw[mmedge] (acc) -- node[above,font=\scriptsize, xshift=0.05cm]{block} (blk);
\draw[mmedge] (v.south west) -- node[right,font=\scriptsize,pos=0.6]{escalate} (esc.east);
\end{mermaidfig}
\vspace{-0.7cm}
\figcaption{Figure 12. Verification before generation. Every proposed robot-patient
command passes a ten-gate assurance suite aligned to fourteen external safety and
robotics standards plus two clinical baselines, with epistemic and aleatory
uncertainty bounds, before the gate surface returns ACCEPT, BLOCK, or ESCALATE
(hand back to human).}

\subsection{Abbreviations}

Table~\ref{tab:abbrev} resolves every abbreviation used across this protocol.

\begin{table}[htbp]
\centering
\footnotesize
\caption{Abbreviations used in this protocol.}
\label{tab:abbrev}
\begin{tabularx}{\textwidth}{L{1.3cm} L{6.8cm} L{1.3cm} L{6.2cm}}
\toprule
\textbf{Abbr.} & \textbf{Definition} & \textbf{Abbr.} & \textbf{Definition} \\
\midrule
AE    & Adverse event                                  & MCP  & Model Context Protocol \\
CTCAE & Common Terminology Criteria for Adverse Events  & MTD  & Maximum tolerated dose \\
DLT   & Dose-limiting toxicity                          & OS   & Overall survival \\
DSMB  & Data and Safety Monitoring Board                & PDAC & Pancreatic ductal adenocarcinoma \\
ECOG  & Eastern Cooperative Oncology Group              & PFS  & Progression-free survival \\
FDA   & U.S. Food and Drug Administration               & PJ   & Pancreaticojejunostomy \\
GCP   & Good Clinical Practice                          & PV   & Portal vein \\
GJ    & Gastrojejunostomy                               & R0   & Margin-negative resection \\
HJ    & Hepaticojejunostomy                             & RP2D & Recommended Phase 2 dose \\
ICH   & International Council for Harmonisation         & SAE  & Serious adverse event \\
IDE   & Investigational Device Exemption                & SMA  & Superior mesenteric artery \\
IND   & Investigational New Drug                        & SMV  & Superior mesenteric vein \\
ISGPS & International Study Group of Pancreatic Surgery & USL  & Unified Safety Level \\
ISM   & Independent Safety Monitor                      & VVUQ & Verification, validation, and uncertainty quantification \\
ITT   & Intention-to-treat                              &      & \\
KRAS  & Kirsten rat sarcoma viral oncogene homolog      &      & \\
LLM   & Large language model                            &      & \\
\bottomrule
\end{tabularx}
\end{table}


